% Generated draft: requires mathematical and source review.
\documentclass[11pt]{article}
\usepackage[margin=25mm]{geometry}
\usepackage{fontspec}
\setmainfont{Latin Modern Roman}
\usepackage{amsmath,amssymb,mathtools,unicode-math}
\setmathfont{Latin Modern Math}
\usepackage{enumitem,needspace,xurl,hyperref}
\setlength{\parindent}{0pt}
\setlength{\parskip}{5pt}
\setlength{\emergencystretch}{3em}
\newcommand{\N}{\mathbb{N}}
\newcommand{\Z}{\mathbb{Z}}
\newcommand{\Q}{\mathbb{Q}}
\newcommand{\R}{\mathbb{R}}
\newcommand{\C}{\mathbb{C}}
\newcommand{\F}{\mathbb{F}}
\newcommand{\A}{\mathbb{A}}
\newcommand{\PP}{\mathbb{P}}
\newcommand{\E}{\mathbb{E}}
\newcommand{\eps}{\varepsilon}
\newcommand{\dd}{\,\mathrm{d}}
\begin{document}
\section*{UnsolvedMath ID: 801046}
\textbf{Source problem code:} catalogue entry
\subsection{Definitions and mathematical statement}
Let \(\mathbb F_2\) be the free group on two generators, with identity \(e\). For a probability measure space \((Y,\nu)\), a probability measure preserving action \(\gamma\) of \(\mathbb F_2\) assigns to each \(g\in\mathbb F_2\) an invertible measurable transformation \(\gamma(g):Y\to Y\) preserving \(\nu\), compatibly with the group operation. Its orbit equivalence relation is
\[
E_\gamma=\{(u,v)\in Y\times Y:\exists g\in\mathbb F_2\; v=\gamma(g)u\}.
\]
Under the usual measure theoretic convention, the action is free if
\[
\nu\bigl(\{u\in Y:\gamma(g)u=u\}\bigr)=0
\qquad\text{for every }g\ne e.
\]
The source does not specify whether freeness instead means pointwise freeness. An action is ergodic if every measurable set invariant modulo null sets under all its transformations has measure zero or one. A particular element \(g\) acts ergodically if its transformation \(\gamma(g)\) is ergodic; explicitly, for every measurable \(A\subseteq Y\),
\[
\nu\bigl(A\mathbin{\triangle}\gamma(g)^{-1}(A)\bigr)=0
\quad\Longrightarrow\quad
\nu(A)\in\{0,1\}.
\]
Here \(\triangle\) denotes symmetric difference.

Let \(\alpha\) be a free, probability measure preserving, ergodic action of \(\mathbb F_2\) on \((X,\mu)\), as in the source. No standardness or nonatomicity hypothesis on this probability space is explicitly supplied.
\begin{enumerate}
\item For every \(x\in\mathbb F_2\setminus\{e\}\), does there exist a free, probability measure preserving action \(\beta\) of \(\mathbb F_2\) on \((X,\mu)\) such that
\[
E_\alpha=E_\beta
\qquad\text{and}\qquad
\beta(x)\text{ is ergodic}?
\]
The action \(\beta\) may depend on \(\alpha\) and on \(x\). The displayed equality of orbit relations is retained from the source; whether equality is required everywhere or only after restriction to a common conull set is unspecified.
\item \textup{(Miller--Tserunyan, as attributed in the supplied record.)} Does there exist a free, probability measure preserving action \(\beta\) of \(\mathbb F_2\) such that
\[
\text{for every }g\in\mathbb F_2\setminus\{e\},
\qquad \beta(g)\text{ is ergodic}?
\]
This asks for one action satisfying the condition simultaneously for every nonidentity element. The source does not explicitly give the underlying probability space for this second action, nor does it repeat \(E_\alpha=E_\beta\). Consequently, it is unresolved whether the intended question requires \(\beta\) on \((X,\mu)\) with \(E_\beta=E_\alpha\), or asks only for an action without that prescribed orbit relation. No such additional requirement is silently imposed here.
\end{enumerate}
\subsection{Short English statement}
Given a free ergodic probability measure preserving action of \(\mathbb F_2\), the first question asks whether its orbit relation admits another free probability measure preserving action making a prescribed nonidentity group element ergodic. The Miller--Tserunyan question asks for a single action making every nonidentity element ergodic. Its intended relationship to the original orbit relation is not explicit in the supplied text.
\subsection{Sources}
\begin{itemize}
\item UnsolvedMath local registry, permanent ID \(801046\), exact source code \texttt{catalogue entry}, title ``Open Problems in Mathematical Logic --- Problem 58.'' The record attributes the problem to Barmpalias, Gao et al., with proposed year \(2026\), and supplies the status label ``open''; that label is not an independently verified mathematical assessment.
\item Supplied source list: Barmpalias, Gao et al., \emph{Open Problems in Mathematical Logic} (2026), Problem 58, \url{https://arxiv.org/abs/2608.26628}. The record reports access on 2026-10-16 and extraction by \texttt{arxiv\_latex}; these are supplied provenance statements, not claims of independent source inspection.
\item The supplied problem text attributes its second question to Miller--Tserunyan. No separate bibliographic reference for that attribution is supplied. The recorded upstream revision is \texttt{d371ee917ac9647fbb06cd6fc529deec6f3dda82}.
\end{itemize}
\end{document}
